% Work-MAT, hodina 54. Podklad: ucebnice Algebraické výrazy, Nasobeni a deleni;
% metodicka inspirace Realisticky.cz 010904; navaznost na 52 a 53.
\documentclass[11pt,a4paper]{article}
% Spolecny styl pro ucitelske pripravy MAT (Work/MAT)
% Verze 1.1, 2026-10-23
\NeedsTeXFormat{LaTeX2e}
\ProvidesPackage{priprava}[2026/10/23 v1.1 Sablona ucitelskych priprav MAT]

% Zakladni typografie a cestina
\RequirePackage[T1]{fontenc}
\RequirePackage[utf8]{inputenc}
\RequirePackage[czech]{babel}
\RequirePackage{lmodern}
\RequirePackage{microtype}
\RequirePackage{amsmath,amssymb}
\RequirePackage{xcolor}
\RequirePackage{enumitem}
\RequirePackage[most]{tcolorbox}
\RequirePackage{tabularx,booktabs,array}
\RequirePackage{fancyhdr}
\RequirePackage{needspace}
\RequirePackage{geometry}
\RequirePackage[hidelinks]{hyperref}

\geometry{a4paper,top=17mm,bottom=18mm,left=23mm,right=23mm}
\definecolor{notegray}{gray}{0.39}
\setlength{\parindent}{0pt}
\setlength{\parskip}{2.5pt plus 1pt minus 0.5pt}
\setlist[itemize]{topsep=2pt,itemsep=1pt,parsep=0pt,leftmargin=1.45em}
\setlist[enumerate]{topsep=2pt,itemsep=2pt,parsep=0pt,leftmargin=1.65em}
\setcounter{secnumdepth}{0}
\renewcommand{\arraystretch}{1.12}

% Zahlavi: vlevo rocnik | globalni tema, vpravo cislo: nazev hodiny
% Zapati zustava zcela prazdne.
\pagestyle{fancy}
\fancyhf{}
\renewcommand{\headrulewidth}{0.35pt}
\renewcommand{\footrulewidth}{0pt}
\setlength{\headheight}{14pt}
\newcommand{\nastavPripravu}[4]{%
  \fancyhead[L]{\small #1\enspace\textbar\enspace #2}%
  \fancyhead[R]{\small #3: #4}%
}

% Minimalisticke cernobile boxy. Kazda definice zvlast.
\tcbset{commonlesson/.style={enhanced,breakable=false,colback=white,colframe=black,
  colbacktitle=white,coltitle=black,boxrule=0.45pt,
  arc=0.8mm,outer arc=0.8mm,left=3mm,right=3mm,top=1.6mm,bottom=1.6mm,
  toptitle=1mm,bottomtitle=0.7mm,fonttitle=\bfseries,
  attach boxed title to top left={xshift=2.5mm,yshift=-1.15mm},
  boxed title style={colback=white,colframe=white,boxrule=0pt,left=1mm,right=1mm}}}
\newtcolorbox{zadaniBox}[1][]{commonlesson,title={Zadání},#1}
\newtcolorbox{prikladBox}[1][]{commonlesson,title={Příklad},#1}
\newtcolorbox{definiceBox}[1][]{commonlesson,title={Definice},#1}

% Spolecny box pro zadani a jeho reseni, bez mezery mezi castmi.
% Pouziti:
%   \begin{prikladReseniBox}
%     Zadani ulohy.
%     \tcblower
%     \begin{align*} ... \end{align*}
%   \end{prikladReseniBox}
% Spodni cast ma jemny sedy podklad a delici linku; lze v ni sazet
% vice radku matematiky. Obe casti maji jen jeden spolecny ramecek.
\tcbset{solutionlesson/.style={
  commonlesson,
  enhanced jigsaw,
  breakable,
  colbacklower=black!3,
  segmentation style={solid,draw=black!25,line width=0.35pt},
  fontlower=\small\color{notegray},
  before lower={\noindent\textbf{Řešení:}\enspace},
  leftlower=3mm,rightlower=3mm
}}
\newtcolorbox{prikladReseniBox}[1][]{solutionlesson,title={Příklad},#1}
\newtcolorbox{zadaniReseniBox}[1][]{solutionlesson,title={Zadání},#1}

% Pomucky pro autora pripravy: sedive poznamky; \reseni pro starsi pripravy
\newcommand{\poznamka}[1]{\par\begingroup\small\color{notegray}\itshape
  \textnormal{\textbf{Poznámka pro učitele.}}\enspace #1\par\endgroup}
\newcommand{\reseni}[1]{\par\begingroup\small\color{notegray}
  \textbf{Očekávané řešení:}\enspace #1\par\endgroup}
\newcommand{\blok}[1]{\needspace{5\baselineskip}\subsection*{#1}}
\newcommand{\vyrok}[1]{\uv{#1}}

\nastavPripravu{1. ročník}{Algebraické výrazy}{54}{Násobení a dělení lomených výrazů}
\begin{document}
\poznamka{45 min: 5 min opakování krácení, 8 min násobení, 12 min dělení a podmínky, 15 min samostatná práce a kontrola, 5 min shrnutí. Rozšíření je volitelné. Nepředpokládej, že nenulovost dělitele automaticky plyne jen z jeho jmenovatele.}
\blok{Násobení: krátíme činitele celých součinů}
Nejdříve zapíšeme podmínky \emph{obou původních zlomků}, potom rozložíme a zkrátíme společné nenulové činitele.
\begin{zadaniBox}[title={Pravidla pro násobení a dělení}]
\[\frac AB\cdot\frac CD=\frac{AC}{BD},\qquad B\ne0,\ D\ne0.\]
\[\frac AB:\frac CD=\frac AB\cdot\frac DC=\frac{AD}{BC},\qquad B\ne0,\ D\ne0,\ C\ne0.\]
Podmínka $C\ne0$ při dělení zaručuje, že \textbf{celý druhý zlomek je nenulový}.
\end{zadaniBox}
\begin{prikladReseniBox}[breakable=false,title={Příklad 1: součin s krácením}]
Uprav $\displaystyle\frac{x^2-4}{x}\cdot\frac{3x}{x+2}$.
\tcblower
Podmínky $x\ne0,-2$. Po rozkladu
\[\frac{(x-2)(x+2)}x\cdot\frac{3x}{x+2}=3(x-2).\]
\end{prikladReseniBox}
\blok{Dělení: dělitel nesmí být nula}
\begin{prikladReseniBox}[breakable=false,title={Příklad 2: zkontroluj tři podmínky}]
Uprav $\displaystyle\frac{x^2-1}{x+2}:\frac{x-1}{x+2}$.
\tcblower
První i druhý zlomek vyžadují $x\ne-2$. Druhý zlomek nesmí být $0$, tedy $x\ne1$.
\[\frac{(x-1)(x+1)}{x+2}\cdot\frac{x+2}{x-1}=x+1,\quad x\ne-2,1.\]
\end{prikladReseniBox}
\begin{prikladReseniBox}[breakable=false,title={Příklad 3: nulový čitatel dělitele}]
Uprav $\displaystyle\frac{x}{x-3}:\frac{2x}{x+1}$.
\tcblower
Z jmenovatelů $x\ne3,-1$ a z nenulovosti druhého zlomku navíc $x\ne0$. Pak
\[\frac{x}{x-3}\cdot\frac{x+1}{2x}=\frac{x+1}{2(x-3)}.\]
\end{prikladReseniBox}
\blok{Samostatná práce}
\begin{zadaniReseniBox}[breakable=false,title={Procvičení: podmínky před každou operací}]
\begin{enumerate}[label=\alph*)]
\item $\frac{3a}{a-1}\cdot\frac{(a-1)^2}{6a}$;
\item $\frac{x^2-4}{x+1}\cdot\frac{x+1}{2-x}$;
\item $\frac{x^2-1}{x}:\frac{x-1}{2x}$;
\item $\frac{x^2-9}{x^2-x}:\frac{x+3}{x}$;
\item $\frac{a+1}{a-2}:\frac{(a+1)^2}{a-2}$.
\end{enumerate}
\tcblower
\begin{enumerate}[label=\alph*)]
\item $\frac{a-1}{2}$, $a\ne0,1$.
\item $-(x+2)$, $x\ne-1,2$.
\item $2(x+1)$, $x\ne0,1$.
\item $\frac{x-3}{x-1}$, $x\ne0,1,-3$.
\item $\frac1{a+1}$, $a\ne2,-1$.
\end{enumerate}
\end{zadaniReseniBox}
\poznamka{Nech vyznačit rozdíl mezi nulovým prvním čitatelem (při jinak splněných podmínkách nevadí) a nulovým druhým zlomkem při dělení (zakázáno).}
\blok{Volitelná rezerva}
\begin{zadaniReseniBox}[breakable=false,title={Rozšíření: srovnej podmínky}]
Pro které hodnoty existuje $\frac{x}{x-1}\cdot\frac{x+2}{x+3}$ a pro které $\frac{x}{x-1}:\frac{x+2}{x+3}$?
\tcblower
Součin: $x\ne1,-3$. Podíl: $x\ne1,-3,-2$, protože při $x=-2$ by byl druhý zlomek nulový.
\end{zadaniReseniBox}
\blok{Domácí procvičení}
\begin{zadaniReseniBox}[breakable=false,title={Domácí úkol: násobení, dělení, podmínky}]
\begin{enumerate}[label=\alph*)]
\item $\frac{x^2-1}{x+3}\cdot\frac{x+3}{x-1}$;
\item $\frac{x}{x-2}:\frac{x+1}{x-2}$;
\item $\frac{a^2-4}{a}\cdot\frac{2a}{a+2}$;
\item $\frac{x^2-9}{x^2-3x}:\frac{x+3}{x}$;
\item $\frac{a+2}{a-1}:\frac{(a+2)^2}{a-1}$;
\item $\frac{x^2-4}{x^2-1}\cdot\frac{x+1}{x+2}$.
\end{enumerate}
\tcblower
\begin{enumerate}[label=\alph*)]
\item $x+1$, $x\ne-3,1$.
\item $\frac{x}{x+1}$, $x\ne2,-1$.
\item $2(a-2)$, $a\ne0,-2$.
\item $1$, $x\ne0,3,-3$; tři zákazy zůstávají i po krácení.
\item $\frac1{a+2}$, $a\ne1,-2$.
\item $\frac{x-2}{x-1}$, $x\ne-2,-1,1$.
\end{enumerate}
\end{zadaniReseniBox}
\end{document}
